\section{Catenary rings}
\label{section-catenary}

\noindent
Compare with Topology, Section \ref{topology-section-catenary-spaces}.

\begin{definition}
\label{definition-catenary}
A ring $R$ is said to be {\it catenary} if for any pair of prime ideals
$\mathfrak p \subset \mathfrak q$, there exists an integer bounding the
lengths of all finite chains of prime ideals
$\mathfrak p = \mathfrak p_0 \subset \mathfrak p_1 \subset \ldots \subset
\mathfrak p_e = \mathfrak q$ and all maximal such chains have the same length.
\end{definition}

\begin{lemma}
\label{lemma-catenary}
A ring $R$ is catenary if and only if the topological space
$\Spec(R)$ is catenary (see
Topology, Definition \ref{topology-definition-catenary}).
\end{lemma}

\begin{proof}
Immediate from the definition and the characterization of
irreducible closed subsets in Lemma \ref{lemma-irreducible}.
\end{proof}

\noindent
In general it is not the case that a finitely generated
$R$-algebra is catenary if $R$ is. Thus we make the following
definition.

\begin{definition}
\label{definition-universally-catenary}
A Noetherian ring $R$ is said to be {\it universally catenary}
if every $R$-algebra of finite type is catenary.
\end{definition}

\noindent
We restrict to Noetherian rings as it is not clear
this definition is the right one for non-Noetherian rings.
By Lemma \ref{lemma-quotient-catenary} to check a Noetherian
ring $R$ is universally catenary, it suffices to check
each polynomial algebra $R[x_1, \ldots, x_n]$ is catenary.

\begin{lemma}
\label{lemma-localization-catenary}
Any localization of a catenary ring is catenary.
Any localization of a Noetherian universally catenary
ring is universally catenary.
\end{lemma}

\begin{proof}
Let $A$ be a ring and let $S \subset A$ be a multiplicative subset.
The description of $\Spec(S^{-1}A)$ in Lemma \ref{lemma-spec-localization}
shows that if $A$ is catenary, then so is $S^{-1}A$. If $S^{-1}A \to C$
is of finite type, then $C = S^{-1}B$ for some finite type ring map
$A \to B$. Hence if $A$ is Noetherian and universally catenary, then
$B$ is catenary and we see that $C$ is catenary too. This proves the lemma.
\end{proof}

\begin{lemma}
\label{lemma-universally-catenary}
Let $A$ be a Noetherian universally catenary ring.
Any $A$-algebra essentially of finite type over $A$
is universally catenary.
\end{lemma}

\begin{proof}
If $B$ is a finite type $A$-algebra, then $B$ is Noetherian
by Lemma \ref{lemma-Noetherian-permanence}. Any finite type
$B$-algebra is a finite type $A$-algebra and hence catenary
by our assumption that $A$ is universally catenary. Thus $B$
is universally catenary. Any localization of $B$ is universally
catenary by Lemma \ref{lemma-localization-catenary} and this
finishes the proof.
\end{proof}

\begin{lemma}
\label{lemma-catenary-check-local}
Let $R$ be a ring. The following are equivalent
\begin{enumerate}
\item $R$ is catenary,
\item $R_\mathfrak p$ is catenary for all prime ideals $\mathfrak p$,
\item $R_\mathfrak m$ is catenary for all maximal ideals $\mathfrak m$.
\end{enumerate}
Assume $R$ is Noetherian. The following are equivalent
\begin{enumerate}
\item $R$ is universally catenary,
\item $R_\mathfrak p$ is universally catenary for all prime ideals
$\mathfrak p$,
\item $R_\mathfrak m$ is universally catenary for all maximal ideals
$\mathfrak m$.
\end{enumerate}
\end{lemma}

\begin{proof}
The implication (1) $\Rightarrow$ (2) follows from
Lemma \ref{lemma-localization-catenary} in both cases.
The implication (2) $\Rightarrow$ (3) is immediate in both cases.
Assume $R_\mathfrak m$ is catenary for all maximal ideals
$\mathfrak m$ of $R$. If $\mathfrak p \subset \mathfrak q$ are primes
in $R$, then choose a maximal ideal $\mathfrak q \subset \mathfrak m$.
Chains of primes ideals between $\mathfrak p$ and $\mathfrak q$ are
in 1-to-1 correspondence with chains of prime ideals between
$\mathfrak pR_\mathfrak m$ and $\mathfrak qR_\mathfrak m$ hence we
see $R$ is catenary. Assume $R$ is Noetherian and $R_\mathfrak m$ is
universally catenary for all maximal ideals $\mathfrak m$ of $R$.
Let $R \to S$ be a finite type ring map. Let $\mathfrak q$ be a prime
ideal of $S$ lying over the prime $\mathfrak p \subset R$.
Choose a maximal ideal $\mathfrak p \subset \mathfrak m$ in $R$.
Then $R_\mathfrak p$ is a localization of $R_\mathfrak m$ hence
universally catenary by Lemma \ref{lemma-localization-catenary}.
Then $S_\mathfrak p$ is catenary as a finite type ring over $R_\mathfrak p$.
Hence $S_\mathfrak q$ is catenary as a localization. Thus $S$ is catenary
by the first case treated above.
\end{proof}

\begin{lemma}
\label{lemma-quotient-catenary}
Any quotient of a catenary ring is catenary.
Any quotient of a Noetherian universally catenary ring is
universally catenary.
\end{lemma}

\begin{proof}
Let $A$ be a ring and let $I \subset A$ be an ideal.
The description of $\Spec(A/I)$ in Lemma \ref{lemma-spec-closed}
shows that if $A$ is catenary, then so is $A/I$.
The second statement is a special case of
Lemma \ref{lemma-universally-catenary}.
\end{proof}

\begin{lemma}
\label{lemma-catenary-check-irreducible}
Let $R$ be a Noetherian ring.
\begin{enumerate}
\item $R$ is catenary if and only if $R/\mathfrak p$ is catenary
for every minimal prime $\mathfrak p$.
\item $R$ is universally catenary if and only if $R/\mathfrak p$ is
universally catenary for every minimal prime $\mathfrak p$.
\end{enumerate}
\end{lemma}

\begin{proof}
If $\mathfrak a \subset \mathfrak b$ is an inclusion of primes of $R$,
then we can find a minimal prime $\mathfrak p \subset \mathfrak a$
and the first assertion is clear. We omit the proof of the second.
\end{proof}

\begin{lemma}
\label{lemma-CM-ring-catenary}
A Noetherian Cohen-Macaulay ring is universally catenary.
More generally, if $R$ is a Noetherian ring and $M$ is
a Cohen-Macaulay $R$-module with $\text{Supp}(M) = \Spec(R)$,
then $R$ is universally catenary.
\end{lemma}

\begin{proof}
Since a polynomial algebra over $R$ is Cohen-Macaulay,
by Lemma \ref{lemma-CM-polynomial-algebra},
it suffices to show that a Cohen-Macaulay ring is
catenary.
Let $R$ be Cohen-Macaulay and $\mathfrak p \subset \mathfrak q$
primes of $R$. By definition $R_{\mathfrak q}$ and
$R_{\mathfrak p}$ are Cohen-Macaulay.
Take a maximal chain of primes
$\mathfrak p = \mathfrak p_0 \subset \mathfrak p_1 \subset
\ldots \subset \mathfrak p_n = \mathfrak q$.
Next choose a maximal chain of primes
$\mathfrak q_0 \subset \mathfrak q_1 \subset \ldots \subset
\mathfrak q_m = \mathfrak p$. By
Lemma \ref{lemma-maximal-chain-CM}
we have $n + m = \dim(R_{\mathfrak q})$. And we have
$m = \dim(R_{\mathfrak p})$ by the same lemma.
Hence $n = \dim(R_{\mathfrak q}) - \dim(R_{\mathfrak p})$
is independent of choices.

\medskip\noindent
To prove the more general statement, argue exactly as above but
using Lemmas \ref{lemma-maximal-CM-polynomial-algebra}
and \ref{lemma-maximal-chain-maximal-CM}.
\end{proof}

\begin{lemma}
\label{lemma-catenary-Noetherian-local}
Let $(A, \mathfrak m)$ be a Noetherian local ring. The following are equivalent
\begin{enumerate}
\item $A$ is catenary, and
\item $\mathfrak p \mapsto \dim(A/\mathfrak p)$ is a dimension function
on $\Spec(A)$.
\end{enumerate}
\end{lemma}

\begin{proof}
If $A$ is catenary, then $\Spec(A)$ has a dimension function $\delta$ by
Topology, Lemma \ref{topology-lemma-locally-dimension-function}
(and Lemma \ref{lemma-catenary}). We may assume $\delta(\mathfrak m) = 0$.
Then we see that
$$
\delta(\mathfrak p) = \text{codim}(V(\mathfrak m), V(\mathfrak p)) =
\dim(A/\mathfrak p)
$$
by Topology, Lemma \ref{topology-lemma-dimension-function-catenary}.
In this way we see that (1) implies (2). The reverse implication
follows from
Topology, Lemma \ref{topology-lemma-dimension-function-catenary}
as well.
\end{proof}














